\clearpage
\begin{thebibliography}{99}
\footnotesize

\bibitem{farkas1902}
J.~Farkas.
\emph{Theorie der einfachen Ungleichungen}.
Journal f\"ur die reine und angewandte Mathematik 124 (1902), 1--27.

\bibitem{schrijver1986}
A.~Schrijver.
\emph{Theory of Linear and Integer Programming}.
Wiley, 1986. (Farkas variants and multiplier sign conventions: \S7.)

\bibitem{krivine1964}
J.-L.~Krivine.
\emph{Anneaux pr\'eordonn\'es}.
Journal d'Analyse Math\'ematique 12 (1964), 307--326.

\bibitem{handelman1988}
D.~Handelman.
\emph{Representing polynomials by positive linear functions on compact
convex polyhedra}.
Pacific Journal of Mathematics 132 (1988), no.~1, 35--62.

\bibitem{tellegen1952}
B.~D.~H.~Tellegen.
\emph{A general network theorem, with applications}.
Philips Research Reports 7 (1952), 259--269.

\bibitem{willems1972}
J.~C.~Willems.
\emph{Dissipative dynamical systems, Part I: General theory}.
Archive for Rational Mechanics and Analysis 45 (1972), 321--351.

\bibitem{moylanhill1978}
P.~J.~Moylan and D.~J.~Hill.
\emph{Stability criteria for large-scale systems}.
IEEE Transactions on Automatic Control 23 (1978), no.~2, 143--149.

\bibitem{vanderschaftjeltsema2014}
A.~van~der~Schaft and D.~Jeltsema.
\emph{Port-Hamiltonian Systems Theory: An Introductory Overview}.
Foundations and Trends in Systems and Control 1 (2014), no.~2--3, 173--378.

\bibitem{jiangteelpraly1994}
Z.-P.~Jiang, A.~R.~Teel, and L.~Praly.
\emph{Small-gain theorem for ISS systems and applications}.
Mathematics of Control, Signals and Systems 7 (1994), 95--120.

\bibitem{jiangmareelswang1996}
Z.-P.~Jiang, I.~M.~Y.~Mareels, and Y.~Wang.
\emph{A Lyapunov formulation of the nonlinear small-gain theorem for
interconnected ISS systems}.
Automatica 32 (1996), no.~8, 1211--1215.

\bibitem{astromfuruta2000}
K.~J.~{\AA}str\"om and K.~Furuta.
\emph{Swinging up a pendulum by energy control}.
Automatica 36 (2000), no.~2, 287--295.

\bibitem{steele2017}
G.~L.~Steele~Jr.
\emph{It's Time for a New Old Language}.
Keynote, ACM SIGPLAN Symposium on Principles and Practice of Parallel
Programming (PPoPP), 2017. (The ``Computer Science Metanotation'' this
book's rule cells are set in.)

\end{thebibliography}
